% =============================================================================
%                    WEEK 2: INTELLIGENCE
% =============================================================================
\section{Week 2: Intelligence}

% -----------------------------------------------------------------------------
%                              2.1 TOPIC SUMMARY
% -----------------------------------------------------------------------------
\subsection{Topic Summary}

\begin{keybox}[Learning Objectives]
By the end of this week, you will be able to:
\begin{enumerate}
  \item Define \textbf{intelligence} and explain why the definition matters.
  \item Distinguish between \textbf{Narrow AI}, \textbf{AGI}, and \textbf{Superintelligence}.
  \item Explain the \textbf{Turing Test} and analyse its interpretations.
  \item Evaluate the \textbf{Chinese Room} argument and its replies.
  \item Understand why intelligence is \textbf{substrate-neutral}.
\end{enumerate}
\end{keybox}

\separator

\subsubsection*{The Big Picture}

This week addresses the foundational question: \textit{What is intelligence?} We explore definitions of intelligence, examine whether machines can be intelligent, and critically analyse two famous thought experiments---the Turing Test and the Chinese Room. These concepts are essential for understanding debates about AI capabilities, consciousness, and the possibility of AGI.

\separator

\subsubsection*{What Examiners Like to Ask}

\begin{itemize}
  \item Define intelligence (Tegmark's definition) and explain Moravec's paradox.
  \item Explain the four quadrants of AI research (Russell \& Norvig).
  \item Describe the Turing Test and its interpretations (sufficient/necessary conditions).
  \item Present the Chinese Room argument and evaluate replies (Systems, Robot, Virtual Mind).
  \item What is substrate neutrality and why does it matter for AI?
\end{itemize}

\bigseparator

% -----------------------------------------------------------------------------
%                 2.2 OVERVIEW + CORE CONCEPTS + EXTENDED THEORY
% -----------------------------------------------------------------------------
\subsection{Overview + Core Concepts + Extended Theory}

\subsubsection*{Roadmap}
\begin{enumerate}
  \item What is intelligence?
  \item Substrate neutrality: information and computation
  \item Four approaches to building AI (Russell \& Norvig)
  \item Types of AI: Narrow AI, AGI, Superintelligence
  \item The Turing Test
  \item Analysis of the Turing Test
  \item The Chinese Room argument
  \item Replies to the Chinese Room
\end{enumerate}

\bigseparator

% --- What is Intelligence? ---
\subsubsection{What is Intelligence?}

\begin{defbox}[Intelligence (Tegmark)]
\textbf{``Intelligence is the ability to accomplish complex goals.''}

This definition:
\begin{itemize}
  \item Subsumes other definitions (logic, understanding, self-awareness, learning).
  \item Allows for \textbf{comparison}: X is more intelligent than Y if X can accomplish all goals Y can, plus at least one more.
  \item Results in a \textbf{partial ordering} (some agents are incomparable).
\end{itemize}
\end{defbox}

\begin{defbox}[Moravec's Paradox]
Tasks that seem \textbf{easy for humans} (image recognition, walking) are \textbf{hard for computers}.

Tasks that seem \textbf{hard for humans} (complex calculations, chess) are \textbf{easy for computers}.

\textbf{Explanation:} Evolution has dedicated massive neural ``hardware'' to sensorimotor tasks, making them feel effortless. Computers lack this specialized circuitry.
\end{defbox}

\begin{tipbox}[Exam Relevance]
Moravec's paradox explains why early AI succeeded at ``hard'' tasks (chess) but struggled with ``easy'' ones (vision). Modern deep learning has begun to address this by learning sensorimotor representations from data.
\end{tipbox}

\bigseparator

% --- Substrate Neutrality ---
\subsubsection{Substrate Neutrality: Information and Computation}

\begin{thmbox}[The Substrate Neutrality Thesis]
Intelligence is ultimately about \textbf{information} and \textbf{computation}, not about the physical material (substrate) that implements them.

\textbf{Implication:} There is no reason in principle why machines (silicon-based) cannot be as intelligent as humans (carbon-based).
\end{thmbox}

\begin{defbox}[Information]
\textbf{Information storage} requires physical components that can exist in multiple \textbf{stable states}.

\textbf{Examples:}
\begin{itemize}
  \item Hard drive: magnetized points (two states $\to$ 1 bit)
  \item RAM: electron positions
  \item Fiber optics: strong/weak laser beams
  \item DNA: nucleotide sequences
  \item Brain: $\sim$10 GB electrically (neuron firing), $\sim$100 TB chemically (synaptic strengths)
\end{itemize}

\textbf{Key point:} Information can be stored in \textbf{any} physical medium with stable states.
\end{defbox}

\begin{defbox}[Computation]
\textbf{Computation} transforms one information state into another by implementing a \textbf{function}.

\textbf{Key result:} NAND gates are \textbf{Turing-complete}---any computable function can be implemented by connecting NAND gates.

\textbf{Implication:} A \textbf{universal intelligent machine} could, given enough time and resources, accomplish any goal as well as any other intelligent entity.
\end{defbox}

\bigseparator

% --- Four Approaches to AI ---
\subsubsection{Four Approaches to Building AI}

Russell and Norvig's classic textbook identifies four paradigms:

\begin{center}
\renewcommand{\arraystretch}{1.4}
\begin{tabular}{@{} l | c c @{}}
\toprule
 & \textbf{Human-Based} & \textbf{Ideal Rationality} \\
\midrule
\textbf{Reasoning-Based} & Think like humans & Think rationally \\
\textbf{Behaviour-Based} & Act like humans & Act rationally \\
\bottomrule
\end{tabular}
\end{center}

\begin{defbox}[Ideal Rational Agent]
A \textbf{perfectly rational agent}:
\begin{enumerate}
  \item Models the environment $E$.
  \item Considers actions $a_1, \ldots, a_n$ and their resulting states $E_1, \ldots, E_n$.
  \item Computes \textbf{expected utility} for each action: $\mathbb{E}[U(a_i)] = P(E_i \mid a_i) \times U(E_i)$.
  \item Chooses the action with highest expected utility.
\end{enumerate}

\textbf{In practice:} Perfect rationality is infeasible. We use \textbf{bounded optimality}---optimal behaviour given resource constraints (time, memory, processing power).
\end{defbox}

\begin{tipbox}[AIXI Framework]
\textbf{AIXI} (Legg \& Hutter) provides a mathematically rigorous definition:

\textit{``Intelligence measures an agent's ability to achieve goals in a \textbf{wide range of environments}.''}

Key insight: The agent doesn't know its environment in advance---it must learn and act simultaneously.

\textbf{Problem:} The optimal AIXI agent is \textbf{not Turing-computable}. Practical AI must approximate this ideal.
\end{tipbox}

\bigseparator

% --- Types of AI ---
\subsubsection{Types of AI: Narrow AI, AGI, Superintelligence}

\begin{center}
\renewcommand{\arraystretch}{1.4}
\begin{tabular}{@{} l p{5cm} p{5cm} @{}}
\toprule
\textbf{Type} & \textbf{Definition} & \textbf{Examples} \\
\midrule
\textbf{Narrow AI} & Purpose-built for \textbf{specific tasks}; often superhuman in that domain & Deep Blue (chess), AlphaGo, GPT-3 \\
\textbf{AGI} & Can accomplish \textbf{any goal} at least as well as humans; possesses common sense, learning, reasoning & Does not yet exist \\
\textbf{Superintelligence} & \textbf{Greatly exceeds} human cognitive performance in virtually all domains & Hypothetical future AI \\
\bottomrule
\end{tabular}
\end{center}

\begin{defbox}[Moravec's Landscape]
Imagine a landscape where:
\begin{itemize}
  \item \textbf{Lowlands} = tasks computers find easy (arithmetic, chess)
  \item \textbf{Peaks} = tasks humans find easy (locomotion, social interaction, creativity)
\end{itemize}

AI capability is like rising water---it has flooded the lowlands and is approaching the peaks. The \textbf{critical sea level} is when machines can perform \textbf{AI design itself}.
\end{defbox}

\begin{thmbox}[The Singularity]
Once machines can \textbf{design better versions of themselves}:
\begin{itemize}
  \item Recursive self-improvement begins.
  \item Machines have advantages over human researchers: speed, memory, data access.
  \item Result: rapid \textbf{intelligence explosion} $\to$ Superintelligence.
\end{itemize}

\textbf{Survey estimates} (Bostrom):
\begin{itemize}
  \item 50\% of researchers expect AGI by 2040.
  \item 10\% expect Superintelligence within 2 years of AGI.
  \item 75\% expect Superintelligence within 30 years of AGI.
\end{itemize}
\end{thmbox}

\bigseparator

% --- The Turing Test ---
\subsubsection{The Turing Test}

\begin{defbox}[The Turing Test (1950)]
\textbf{Original formulation (Imitation Game):}
\begin{itemize}
  \item An interrogator communicates (via text) with a human and a machine in separate rooms.
  \item The interrogator tries to determine which is which.
  \item The machine tries to fool the interrogator; the human tries to help identify the machine.
\end{itemize}

\textbf{Turing's prediction:} In 50 years, a computer could fool an average interrogator 30\% of the time after 5 minutes.

\textbf{Alternative formulation:} If the judge cannot do better than 50/50, the machine \textbf{passes} the test.
\end{defbox}

\begin{tipbox}[Historical Note]
Turing anticipated \textbf{machine learning}: ``Instead of trying to produce a programme to simulate the adult mind, why not\ldots simulate the child's? If\ldots subjected to an appropriate course of education one would obtain the adult brain.''

Descartes (1669) prefigured the test: machines could never ``arrange speech in various ways, in order to reply appropriately to everything that may be said in its presence.''
\end{tipbox}

\bigseparator

% --- Analysis of the Turing Test ---
\subsubsection{Analysis of the Turing Test}

\begin{defbox}[Logical Interpretations]
The Turing Test can be interpreted as providing:

\textbf{1. Logically sufficient conditions:}
\[
\text{Pass TT} \Rightarrow \text{Intelligent}
\]
(TT-passers are a \textbf{subset} of intelligent beings. Failing TT doesn't prove non-intelligence.)

\textbf{2. Logically necessary conditions:}
\[
\text{Intelligent} \Rightarrow \text{Pass TT}
\]
(Intelligent beings are a \textbf{subset} of TT-passers. Failing TT proves non-intelligence.)

\textbf{3. Both (biconditional):} Pass TT $\Leftrightarrow$ Intelligent. (Rarely claimed.)

\textbf{4. Probabilistic:} Passing TT provides \textbf{evidence} for intelligence (Turing's own view).
\end{defbox}

\begin{exbox}[Objections to the Turing Test]
\textbf{1. The Chauvinistic Objection} (attacks necessary interpretation):
\begin{itemize}
  \item Intelligent aliens might fail because they don't share human linguistic conventions.
  \item An entity could achieve complex goals in diverse environments without passing TT.
\end{itemize}

\textbf{2. The Blockhead Objection} (attacks sufficient interpretation):
\begin{itemize}
  \item Imagine a creature with a \textbf{lookup table} containing responses to every possible input.
  \item Blockhead passes TT but uses \textbf{brute force}---no real intelligence or understanding.
  \item \textbf{Counter:} Is Blockhead logically/nomically possible? Is brute-force processing really non-intelligent?
\end{itemize}

\textbf{3. Too Hard:} Nothing without a human-like brain could pass.

\textbf{4. Too Narrow:} Success might come from tricks, not genuine cognition.

\textbf{5. Too Easy:} More demanding tests needed (e.g., \textbf{Lovelace Test}---creativity/originality).
\end{exbox}

\begin{tipbox}[Exam Tip]
When discussing TT interpretations, use set-theoretic reasoning:
\begin{itemize}
  \item Sufficient: $\{\text{TT-passers}\} \subseteq \{\text{Intelligent beings}\}$
  \item Necessary: $\{\text{Intelligent beings}\} \subseteq \{\text{TT-passers}\}$
\end{itemize}
Then apply this to analyse objections.
\end{tipbox}

\bigseparator

% --- The Chinese Room ---
\subsubsection{The Chinese Room Argument (Searle, 1980)}

\begin{thmbox}[The Chinese Room]
\textbf{Setup:}
\begin{itemize}
  \item A person who knows no Chinese is locked in a room.
  \item The room contains boxes of Chinese symbols (database) and a rulebook (program).
  \item People outside send in Chinese questions (input).
  \item The person follows rules to manipulate symbols and produces Chinese answers (output).
  \item The person passes the Turing Test for Chinese---but understands nothing.
\end{itemize}

\textbf{Argument structure:}
\begin{enumerate}
  \item The symbol-manipulating program $P$ is possible (at least conceivably).
  \item $P$ does not possess \textbf{understanding} $\to$ $P$ is not intelligent.
  \item Any digital computer $C$ is equivalent to $P$ in relevant respects.
  \item \textbf{Conclusion:} No digital computer can be intelligent.
\end{enumerate}
\end{thmbox}

\begin{defbox}[Key Distinction: Syntax vs.\ Semantics]
\begin{itemize}
  \item \textbf{Syntax:} Formal manipulation of symbols according to rules.
  \item \textbf{Semantics:} Meaning---what symbols \textbf{represent}.
\end{itemize}

\textbf{Searle's claim:} Computers can only do syntax. Understanding requires semantics. Therefore, computers cannot truly understand---and understanding is necessary for intelligence.

\textbf{Challenge:} Is understanding really necessary for intelligence (per Tegmark's definition)?
\end{defbox}

\bigseparator

% --- Replies to the Chinese Room ---
\subsubsection{Replies to the Chinese Room}

\begin{center}
\renewcommand{\arraystretch}{1.4}
\begin{tabular}{@{} l p{9cm} @{}}
\toprule
\textbf{Reply} & \textbf{Argument} \\
\midrule
\textbf{Systems Reply} & The \textbf{man} doesn't understand, but the \textbf{whole system} (man + database + rulebook) does. The man is just the CPU. \\
\textbf{Virtual Mind Reply} & The system creates a \textbf{virtual entity} (like Siri or a video game character) that understands, distinct from the physical components. \\
\textbf{Robot Reply} & Understanding requires \textbf{embodiment}---interacting with the world. A robot that sees, touches, and experiences could ground symbols in meaning. \\
\textbf{Impossibility Reply} (Dennett) & Such a program would require $\sim$100 billion lines of code. Running it by hand would take lifetimes. It's not \textbf{nomically possible}. \\
\bottomrule
\end{tabular}
\end{center}

\begin{defbox}[Nomic vs.\ Logical Possibility]
\begin{itemize}
  \item \textbf{Logical possibility:} No inherent contradiction (can be conceived).
  \item \textbf{Nomic possibility:} Possible given the \textbf{laws of physics}.
\end{itemize}

Dennett argues the Chinese Room is \textbf{logically} possible but not \textbf{nomically} possible---so it's not a valid counterexample.
\end{defbox}

\begin{tipbox}[Exam Tip]
The Chinese Room hinges on \textbf{Premise 2}: understanding is necessary for intelligence. If you accept Tegmark's goal-based definition, you can reject this premise---a system could be intelligent (achieve complex goals) without ``true understanding.''
\end{tipbox}

\bigseparator

% --- Summary ---
\subsubsection{Summary: Key Takeaways}

\begin{thmbox}[Week 2 Key Points]
\begin{enumerate}
  \item \textbf{Intelligence} = ability to accomplish complex goals (Tegmark).
  \item \textbf{Moravec's paradox:} ``Easy'' human tasks are computationally hard.
  \item \textbf{Substrate neutrality:} Intelligence depends on information/computation, not physical material.
  \item \textbf{Four AI paradigms:} Think/act like humans vs.\ think/act rationally.
  \item \textbf{Narrow AI} excels at specific tasks; \textbf{AGI} matches human breadth; \textbf{Superintelligence} exceeds humans in all domains.
  \item \textbf{Turing Test:} Linguistic indistinguishability as a test for intelligence.
  \item TT can be interpreted as sufficient, necessary, both, or probabilistic.
  \item \textbf{Chinese Room:} Syntax alone doesn't yield understanding.
  \item Replies: Systems, Virtual Mind, Robot (embodiment), Impossibility.
\end{enumerate}
\end{thmbox}

\begin{defbox}[Key Terms to Remember]
\begin{itemize}
  \item \textbf{Narrow AI / AGI / Superintelligence:} Levels of AI capability.
  \item \textbf{Substrate neutrality:} Intelligence is hardware-independent.
  \item \textbf{Turing-complete:} Can compute any computable function.
  \item \textbf{Expected utility:} $P(\text{outcome}) \times U(\text{outcome})$.
  \item \textbf{Bounded optimality:} Best behaviour given resource limits.
  \item \textbf{Turing Test:} Behavioural test for machine intelligence.
  \item \textbf{Logically sufficient/necessary:} Set-theoretic interpretations.
  \item \textbf{Chinese Room:} Thought experiment against computer understanding.
  \item \textbf{Syntax vs.\ Semantics:} Symbol manipulation vs.\ meaning.
  \item \textbf{Nomic possibility:} Possible given physical laws.
  \item \textbf{Embodiment:} Understanding requires physical interaction with the world.
\end{itemize}
\end{defbox}

